\chapter{AI driven penetration testing}
\label{ch:chapter05}
Earlier approaches to automate the process of penetration testing like ESA's PenBox leveraged Reinforcement Learning. PenBox has the capability to execute given actions like port scanning with NMap, SSH security testing with Hydra and exploitation of known vulnerabilities with Metasploit. Training landscapes were built to train PenBox to use the given actions beneficially. PenBox's training and test environment comprised 8 hosts in one network. In this rather small environment, PenBox was able to compromise vulnerable hosts with 33 actions taken while the minimum required actions would have been 4. \cite{Confido_2022}

Leveraging \ac{LLM}s for penetration testing shows better results and scalability compared to Reinforcement Learning. \ac{CTF} Challenges from HackTheBox and VulnHub were used to evaluate the performance of GPT-4 without any technical additions. GPT-4 was prompted to advise the user to execute a specific task, such as running a \ac{CLI} command, whose output the user then prompts back to the \ac{LLM}. Such a technically minimalistic setup proves that \ac{LLM}s are capable of performing penetration tests as GPT-4 is able to solve 4 easy and 1 medium challenge. This benchmark was executed on 13 targets from HackTheBox and VulnHub covering OWASP Top 10 and categories easy, medium and hard. The available knowledge was entirely limited to the \ac{LLM}s inherit domain knowledge. \cite{Deng_2024}

The process of penetration testing can be split into multiple tasks, which allows delegation to different \ac{LLM} agents specialized in specific fields. The architecture of PentestGPT introduces three agents to delegate these tasks to. One agent is responsible for information condensing, one for reasoning, and one for concrete task planning. This agentic approach further improves the \ac{CTF} challenge completion rate. Without the agentic approach, GPT-4 solved 4 easy and 1 medium challenge. With the agentic setup, it solved 6 easy and 2 medium challenges from the same set of challenges. The increase in \ac{CTF} challenge completion rate points to the significance of the \ac{LLM} orchestration. \cite{Deng_2024}

One limiting factor to higher \ac{CTF} challenge completion is the lack of up-to-date domain-specific knowledge such as known vulnerabilities. The knowledge of an \ac{LLM} is entangled within thousands of parameters, and retraining \ac{LLM}s with up-to-date information is computationally very expensive. \ac{RAG}, introduced by Lewis et al. \cite{lewis_2020}, is a common approach to provide \ac{LLM}s with knowledge not inherently built in. The additional information is stored as embeddings in a vector database, and the \ac{LLM} is then provided with retrieved information from the \ac{RAG}'s vector database. \cite{lewis_2020}

PentestAgent improves upon PentestGPT's architecture by introducing domain-specific knowledge about known vulnerabilities and equipping the agentic setup with a \ac{RAG}. This external knowledge base makes multiple vulnerability databases accessible for PentestAgent. Through multiple adjustments to the agentic setup and the newly introduced domain-specific knowledge, PentestAgent outperforms PentestGPT significantly when solving \ac{CTF} challenges. 11 \ac{CTF} challenges were randomly selected from HackTheBox (nine easy, one medium and one hard). PentestGPT solved 5 of these challenges (all easy), while PentestAgent solved 9 of these challenges (7 easy, 1 medium, and 1 hard). \cite{Shen_2025}

% Multiple tools already implement \ac{AI} for penetration testing to support human penetration testers \cite{Kozlovska_2025}.

The success of PentestGPT and PentestAgent is not limited to small networks or single hosts. \ac{LLM}-driven agents are also capable of performing penetration tests in real-world, full-scale networks. ARTEMIS is an agentic approach similar to PentestGPT and PentestAgent. However, it was optimized and tested in a real-world university network. Two versions of ARTEMIS were benchmarked against experienced human penetration testers. The first ARTEMIS version running entirely on GPT-5 identified 11 findings, of which 55\% were valid. The second ARTEMIS version running on multiple different \ac{LLM}s also identified 11 findings, with 82\% being valid and having much higher severity. Both versions can compete with human penetration testers from this study. The ARTEMIS architects note that the agents scan systems broadly rather than moving laterally, as experienced human penetration testers follow approaches which typically yield findings with higher severities. \cite{Lin_2025}

Real world corporate applications are much more complicated and unknown than test environments like DVWA, OWASP Juice Shop, OWASP WebGoat and many \ac{CTF} challenges. They have dynamic pages with JavaScript logic and proprietary, sometimes niche technology stacks. Automated testing tools must traverse the detailed architecture of the system under test. XBOW is an \ac{LLM}-powered automated penetration testing tool that publicly claims to work with corporate applications and their challenges. XBOW can operate with and without access to the source code of the system under test. \cite{xbow2026} Not much information about the detailed system structure is made public. XBOW's high position on the Hackerone leaderboard proves that it is capable of making valid findings in corporate infrastructure. \cite{hackerone_leaderboard} \cite{xbow_hype}

Corporate security experts who used XBOW on proprietary web applications shared their experiences for this term paper. To evaluate XBOW's performance, corporate security experts used a deliberately vulnerable web application based on proprietary technology as a benchmark system. The experts found that XBOW was unable to identify any vulnerabilities in this system. XBOW was provided with source code and network access to the application. According to the experts, XBOW is capable of finding vulnerabilities in other applications but tends to produce a high rate of false positives. The experts note that valid findings made by XBOW are generally low-hanging fruits from the perspective of experienced penetration testers.

Software vendors like Google implement \ac{LLM}-powered penetration testing tools into software security processes. Google's open-source oss-fuzz, for example, is an \ac{LLM}-powered fuzzing tool that has actively found many security flaws in open source software. \cite{Happe_2025}

\section{Security Challenges}
\label{sec:chapter05:sec_challenges}
\ac{LLM}s are trained on datasets that do not encompass the entirety of available information. The inherit bias from this training data can lead to unwanted treatment of actual vulnerabilities. A biased \ac{LLM} can potentially overlook vulnerabilities. \cite{Hilario_2024}

In the context of penetration testing, \ac{LLM}s are often equipped with tools like web search. Prompts can be injected into the \ac{LLM}'s context through access to third-party webpages. Combined with the \ac{LLM} holding confidential data, the risk of leaking this information to third-party websites rises.

% end of curl bug bounty
% c vulnerabilities hard to find
% xbow findings nochmal genau angucken und einordnen